\begin{proofbox}
\paragraph{最短通用证明：中值定理}
任取 $x_1<x_2$，$x_1,x_2\in I$。区间性保证 $[x_1,x_2]\subset I$，连续与可导条件保证存在 $\xi\in(x_1,x_2)$，使
\[
f(x_2)-f(x_1)=f'(\xi)(x_2-x_1).
\]
因为 $x_2-x_1>0$，右端与 $f'(\xi)$ 同号。这一次计算同时证明严格递增、严格递减、不减和不增四个充分性结论。

反过来，若 $f$ 不减，对内部点 $x$，只要 $h\ne0$ 且 $x+h\in I$，就有
\[
\frac{f(x+h)-f(x)}{h}\ge0.
\]
令 $h\to0$ 得 $f'(x)\ge0$。不增的情形将不等号反向。

最后，若 $f'\ge0$ 却不严格递增，则由不减性，存在 $x_1<x_2$ 使 $f(x_1)=f(x_2)$。对 $x\in[x_1,x_2]$，有 $f(x_1)\le f(x)\le f(x_2)$，所以 $f$ 在该段恒定。
因此，排除常值子区间就得到严格递增；有限个零导数点不可能覆盖一个非退化区间。

\textbf{知识门槛说明：}一般导数判定定理需要中值定理作为依据；具体多项式题常可直接作差，不必每次重复引用该定理。
\end{proofbox}
